%==============================================================================
% Benchmark Construction Detail --- appendix.
% Authored by the design agent. Gathers content relocated out of the body to keep
% the Design section at ~1.25pp:
%
%   FROM sections/design.tex (body):
%     * the full Helix company-skeleton enumeration (founders, per-year arcs,
%       documentation-discipline eras, product lines, tooling-replacement timeline)
%     * the expanded "why synthetic" control + privacy reasons
%     * the ground-truth / rubric field schemas      --- carries \label{sec:design:rubric}
%     * the persona library + document-genre taxonomy + artifact counts
%                                                     --- carries \label{sec:design:personas}
%
%   FROM the gutted sections/generation.tex (handed to design):
%     * the Stage~A--F generation pipeline
%     * the genre taxonomy
%     * the documentation-discipline eras
%     * the full circularity structural-defenses enumeration
%
% This file is \input from main.tex's appendix block. It does NOT re-issue
% \appendix. It defines \label{app:construction}. The body refs in design.tex
% (\S\ref{app:construction}) and the pointer in sections/generation.tex resolve
% here. Stage~E references the C1--C6 taxonomy, which lives in the Design body, so
% it \ref's tab:categories / sec:design:taxonomy normally.
%
% This file is now ALSO the physical home of the per-category question-count
% table \label{tab:design:counts} (relocated out of sections/design.tex during
% the body-shortening pass). It sits beside sec:design:personas; design.tex now
% only \ref's it. Do not duplicate this \label back into design.tex.
%==============================================================================

\section{Benchmark Construction Detail}
\label{app:construction}

This appendix collects the construction detail relocated from the Design section
(Section~\ref{sec:design}): the full Helix company skeleton, the remaining
arguments for a synthetic corpus, the ground-truth and rubric field schemas, the
persona library and document-genre taxonomy with artifact counts, and the
stage-by-stage generation pipeline.

\subsection{The Helix Company Skeleton}
\label{app:construction:helix}

The seed organization, \textbf{Helix Logistics}, is instantiated with a coherent
multi-year narrative arc: founded in 2020 by Maya Patel and Daniel Kim, shipping
the Helix API in 2021, a Series~A and a business-to-consumer to mid-market pivot
in 2022, a 2023 customer escalation accompanied by a fifteen-percent layoff and a
change of COO, the Helix AI Optimizer in 2024, a documentation overhaul in 2025,
and a planned Series~B in 2026. A fixed company skeleton pins the founding year,
founders, per-year arcs, documentation-discipline eras, product lines,
customer-arc count, and a tooling timeline with explicit replacements (for
example, Salesforce replaces HubSpot in 2022-Q1 and Confluence replaces Notion in
2023-Q1). These replacements are not decoration: they seed the temporal and
supersession structure that the hard question categories then probe.

\subsection{Why Synthetic: Control and Privacy}
\label{app:construction:why-synthetic}

The body (Section~\ref{sec:design:why-synthetic}) states the decisive reason for a
synthetic corpus---no overlap with any system's pretraining---in full. Two further
reasons complete the rationale. First, \textbf{control}: the texture is
realistic---drafts, contradictions, half-finished threads, decisions that
change---while the ground truth is fully known, because answers are projected from
the source-of-truth graph rather than adjudicated after the fact from messy free
text. A real, scraped corpus would force a labor-intensive and error-prone
annotation effort and would still leave the gold answers contestable; here the
answer key is constructed before the evidence, and the evidence is generated to
license it. Second, \textbf{privacy}: no real company's confidential
communications---emails, Slack messages, contracts, incident reports---are
exposed, which both removes a serious release hazard and obviates the personally
identifiable information scrubbing that any real organizational corpus would
require. The cost is the well-known generalization gap between synthetic and real
corpora, which we address directly in our limitations: synthetic texture is a
model of organizational communication, not a sample of it, and we make no claim
that absolute scores transfer unchanged to a specific real company.

\subsection{Ground Truth and Rubric Structure}
\label{sec:design:rubric}

Because the corpus is the \emph{evidence} and the source-of-truth graph is the
\emph{answer key}, every ground-truth answer is a deterministic projection of the
graph rather than free text recalled by a model. Each question record carries an
identifier, a category, a natural-language question \texttt{text}, three
paraphrases (so that scoring is robust to phrasing), a structured
\texttt{ground\_truth\_answer}, weighted \texttt{rubric\_subpoints}, a list of
\texttt{evidence\_artefact\_ids} naming the specific corpus artifacts that
license the answer, and metadata---including a \texttt{leak\_score} that measures
how much the surface question text gives the answer away, plus dated audit and
provenance notes recorded across validation waves.

% {\sloppy} + a local \emergencystretch let TeX relax inter-word spacing across
% this paragraph's long unbreakable \texttt identifiers (the canonical, \ref'd
% field-schema names), clearing the residual overfull \hbox warnings without
% altering any identifier. The identifiers stay verbatim by design.
{\sloppy\emergencystretch=3em
The \texttt{ground\_truth\_answer} is a set of category-specific structured
fields, not a single string, which is precisely what makes it a projection of
the graph. C1 records \texttt{current\_value}, \texttt{prior\_value},
\texttt{change\_date}, \texttt{change\_reason}, and a \texttt{scope\_note}. C2
records the \texttt{decision}, its \texttt{deciders}, the \texttt{decision\_date},
the \texttt{alternatives}, and the \texttt{deciding\_factor}. C3 separates
\texttt{original\_value} and \texttt{corrected\_value} from
\texttt{original\_recorded\_at}, \texttt{correction\_recorded\_at}, and
\texttt{true\_occurred\_at}, making the record-time/world-time distinction
explicit. C4 records an \texttt{as\_of\_date} together with \texttt{asof\_facts},
\texttt{changed\_facts}, and \texttt{replacements}. C5 records a \texttt{claim},
an \texttt{evidence\_summary}, \texttt{evidence\_types},
\texttt{testimony\_attribution}, and \texttt{inferential\_steps}. C6 records
\texttt{party\_a}, \texttt{party\_b}, each party's claim and date, a
\texttt{resolution\_status}, and \texttt{resolution\_details}. The large-tier
emergent-pattern category (out of scope here) records a ground-truth answer
alongside its \texttt{facets}, a \texttt{temporal\_qualifier}, and a
\texttt{formalization} of the de-facto rule.
\par}

Scoring is \textbf{rubric-based facet coverage}, not exact match. Each question
carries weighted sub-points---C1 has four sub-points at $0.25$ each; C3, C4, C5,
and C6 typically carry three sub-points at roughly $0.333$ each. (Large-tier
emergent-pattern questions, out of scope here, carry seven to eight facet
sub-points with explicit per-facet \texttt{weight} and \texttt{evidence\_ids}.)
Each sub-point pairs a positive
\texttt{criterion} (what must be asserted, with paraphrase accepted) with an
explicit \texttt{fail\_criterion} (what voids credit). For example, a C2
sub-point requires naming all of the relevant deciders and fails only if none are
mentioned; a C5 sub-point awards full credit for listing all three evidence items
with the correct primary/secondary roles, half credit for two, and zero below
two. This facet-coverage structure lets multiple correct phrasings earn full
credit while omissions are penalized deterministically. The
\texttt{evidence\_artefact\_ids} field ties each answer to its licensing
artifacts---one C1 question, for instance, is grounded in six artifacts spanning
two distinct events---which enables per-question evidence audits and makes the
basis for every gold answer inspectable.

\subsection{Persona Library and Document Genres}
\label{sec:design:personas}

A multi-author corpus is only meaningful if its authors have distinct, consistent
voices, so \bench{} draws from a persona library of $72$ named personas. Each
persona record carries a stable identifier and display name, pronouns, join and
(where applicable) departure years, a dated \texttt{role\_history} so that
seniority and reporting lines move over time, \texttt{voice\_markers} (hedging
frequency, emoji use, formality, typical message length, whether the person
writes long documents), three to five \texttt{signature\_phrases}, a
documentation-discipline trait, a per-channel communication mix, and
domain-expertise tags, plus \texttt{is\_long\_tail} and \texttt{is\_founder}
flags. This gives each author a recognizable voice across the corpus---Maya
Patel, the CEO, writes terse and data-driven with low documentation discipline
(``Numbers don't lie''), while Daniel Kim, the CTO, writes long-form narrative
with high documentation discipline (``Let me walk you through'')---and it is what
makes the testimony-attribution facets of C5 and the party-attribution facets of
C6 meaningful: the same fact reads differently depending on who said it and where
in their role history they said it. The \texttt{is\_long\_tail} flag seeds the
sparse, rarely mentioned actors that defeat naive retrieval, which is exactly the
population an organizational memory system must not lose.

The corpus is organized as event folders, each holding one or more artifacts that
report the same event from different angles. Every artifact begins with a
metadata block recording its slot ID, event ID, genre, role (primary versus
secondary witness to the event), author (a persona), and testimony type
(\texttt{direct} versus \texttt{inferred}). The genres mirror where work actually
accumulates; in the medium tier the artifact counts are: Slack threads ($115$),
email threads ($109$), meeting notes ($102$), meeting transcripts ($65$),
incident reports ($24$), sales/CRM entries ($11$), Notion documents ($9$),
postmortems ($7$), and one architecture decision record; the small tier spans the
same genres and additionally includes retrospectives. The texture is realistic by
construction: meeting transcripts include timestamped turns with
\texttt{[crosstalk]} and \texttt{[inaudible]} markers and natural disfluencies,
email threads carry full headers and signature blocks, and the same decision is
deliberately surfaced in more than one channel---a transcript in which two
executives talk through a pivot to high-value accounts, followed by an email
thread that confirms it for the sales team---so that an answer is licensed by
\emph{cross-channel agreement} rather than a single canonical source. At medium
scale the corpus comprises $443$ artifacts across $157$ events, split primary
$152$ / secondary $291$ and direct $289$ / inferred $154$: most facts must be
triangulated from secondary, partly inferential witnesses rather than read off a
single record. The small tier holds the same genre mix at reduced scale: $121$
corpus artifacts across $17$ events. This composition is unchanged from earlier
builds: the corpus still physically contains the emergent-evidence background
documents (retained as benign background noise), and only the emergent
\emph{questions} were removed when that category was deferred to the large tier;
the corpus itself was not regenerated.

Table~\ref{tab:design:counts} gives the per-category question counts over this
corpus at both tiers; the Design section (Section~\ref{sec:design:taxonomy})
points here for the full breakdown.

\begin{table}[t]
  \centering
  \caption{\textbf{Question counts per category, by tier.} The emergent-pattern
    category is evaluated only at the large tier and above and is out of scope
    for the small and medium tiers. Categories are defined in
    Table~\ref{tab:categories} and Section~\ref{sec:design:taxonomy}.}
  \label{tab:design:counts}
  \begin{tabular}{l l c c}
    \toprule
    Code & Category & Small ($n$) & Medium ($n$) \\
    \midrule
    C1           & Supersession        & 1  & 15 \\
    C2           & Decision provenance & 2  & 15 \\
    C3           & Bi-temporal         & 3  & 7  \\
    C4           & Audit replay        & 1  & 15 \\
    C5           & Justification chain & 3  & 15 \\
    C6           & Contradiction       & 1  & 6  \\
    \midrule
    \multicolumn{2}{l}{\textbf{Total}} & \textbf{11} & \textbf{73} \\
    \bottomrule
  \end{tabular}
\end{table}

\subsection{Generation Pipeline}
\label{app:construction:pipeline}

The Helix corpus is entirely synthetic: no human authored any artifact in the
dataset. As summarized in the body (Section~\ref{sec:design:generation}), every
artifact is produced by the parameterized, template-driven \texttt{helix-corpus}
pipeline driving a single \emph{local} open-weight model (\texttt{gpt-oss:20b},
served via Ollama). The only human input is the choice of scale parameters that
seed the fictional company---its size, year range, and team topology---after which
generation proceeds through six ordered stages.

\paragraph{Stages.}
The pipeline moves from an abstract company description to a fully validated
question set, projecting ground truth deterministically from a structured
graph rather than recalling it from free text.
%
\emph{Stage~A (canon)} instantiates the fictional company, Helix Logistics:
founders, multi-year growth arcs, a persona registry of several dozen named
employees with titles and tenure dates, customer arcs, documentation-discipline
eras, and a tooling timeline with explicit replacements. The output is structured Markdown and YAML, plus a fixed
\texttt{skeleton.json} that seeds the temporal and supersession structure used
downstream.
%
\emph{Stage~B (source-of-truth graph)} converts the canon into a ground-truth
event log and bi-temporal relation graph: each event is keyed to a date range,
a category, a set of participant persona IDs, and a canonical fact or decision,
and relations encode supersessions, corrections, and contradictions. Crucially,
the answers to every downstream question are projected deterministically from
this graph; the model is never asked to invent ground truth, only to read it
off a structured record.
%
\emph{Stage~C (corpus)} generates a cluster of two to three realistic artifacts
per event---Slack threads, meeting notes, email chains, retrospectives, ADRs,
and Notion-style documents---with genre assigned per event from a fixed
taxonomy. Each artifact carries a structured header recording its slot ID,
event ID, author, date, and testimony type (\texttt{direct} versus
\texttt{inferred}).
%
\emph{Stage~D (cross-references and noise)} adds cross-artifact references and
injects realistic distractors---drafts, half-finished threads, contradictions
among secondary witnesses, and atmosphere artifacts---recording every injected
distractor in a noise manifest so the corpus stays auditable.
%
\emph{Stage~E (questions)} generates the natural-language question texts and
paraphrases, with answers computed from the Stage~B graph rather than authored
by the model; questions span the six structured categories (C1--C6; defined in
Table~\ref{tab:categories} and Section~\ref{sec:design:taxonomy}) plus an
EMERGENT tier, and each carries its canonical answer, evidence artifact IDs, and
a weighted rubric.
%
\emph{Stage~F (validation)} runs a self-consistency pass over the generated
output, checking answerability, rubric coherence, and distractor integrity;
questions that pass become the shipped \texttt{benchmark\_v0.0.jsonl}.

\paragraph{Why this design bounds, but does not eliminate, circularity.}
Single-model end-to-end generation is circular by construction: the same model
writes the artifacts, phrases the questions, and performs the self-consistency
check. Structural defenses---closed-enum event types, graph-projected answers,
multi-pass critique, and automated sentinels for persona consistency and
anachronism---reduce this circularity but do not remove it. Following
\citet{gebru2021datasheets} we document it rather than obscure it, and a
cross-model validation pass (a frontier model used as an independent validator,
not as the generator) is planned. The three-wave manual audit and the six defect
classes it corrected are enumerated in Appendix~\ref{app:annotation} and
Section~\ref{sec:generation-corrections}.
